% Part: model-theory
% Chapter: basics
% Section: dlo

\documentclass[../../../include/open-logic-section]{subfiles}

\begin{document}

\olfileid{mod}{bas}{dlo}
\section{الترتيبات الخطية الكثيفة}

\begin{defn}
  \emph{الترتيب الخطي الكثيف بلا طرفين} هو !!a{structure}~$\Struct{M}$
  للـ!!{language} التي تحتوي !!{predicate} واحدًا ذا موضعين~$<$ يحقق
  الجمل الآتية:
  \begin{enumerate}
  \item $\lforall[x][\lnot x < x]$؛
  \item $\lforall[x][\lforall[y][\lforall[z][(x < y \lif (y < z \lif x
    <z ))]]]$؛
  \item $\lforall[x][\lforall[y][(x< y \lor \eq[x][y] \lor y < x)]]$؛
  \item $\lforall[x][\lexists[y][x < y]]$؛
  \item $\lforall[x][\lexists[y][y < x]]$؛
  \item $\lforall[x][\lforall[y][(x < y \lif \lexists[z][(x < z \land
        z < y))]]]$.
 \end{enumerate}
\end{defn}

\begin{thm}\ollabel{thm:cantorQ}
  كل ترتيبين خطيين كثيفين !!{enumerable} بلا طرفين متشاكلان.
\end{thm}

\begin{proof}
  لتكن $\Struct{M_1}$ و$\Struct{M_2}$ ترتيبين خطيين كثيفين
  !!{enumerable} بلا طرفين، حيث ${<_1} = \Assign{<}{M_1}$ و
  ${<_2} = \Assign{<}{M_2}$، ولتكن $\PIso{I}$ مجموعة جميع التشاكلات
  الجزئية بينهما. ليست $\PIso{I}$ خالية، إذ إن
  $\emptyset \in \PIso{I}$ على الأقل. نبين أن $\PIso{I}$ تحقق خاصية
  الذهاب والإياب. وعندئذ $\Struct{M_1} \iso[p] \Struct{M_2}$، وتتبع
  المبرهنة من \olref[pis]{thm:p-isom1}.

  لإثبات أن $\PIso{I}$ تحقق خاصية الذهاب، خذ $p \in \PIso{I}$ و
  $a \in \Domain{M_1}$. إذا كانت $p$ الدالة الخالية، فاختر أي
  $b \in \Domain{M_2}$ وخذ $q = \{\langle a,b\rangle\}$. وإذا كان
  $a$ في مجال $p$، فخذ $q=p$. ويبقى أن تكون $p$ غير خالية وألا يكون
  $a$ في مجالها. اكتب عندئذ $p(a_i)=b_i$، حيث $i=1$، \dots،~$n$،
  وافترض من غير إخلال بالعموم أن $a_1 <_1 a_2 <_1 \cdots <_1 a_n$.
  أوجد $b \in \Domain{M_2}$ على النحو الآتي:
  \begin{enumerate}
  \item إذا كان $a <_1 a_1$، فليكن $b \in \Domain{M_2}$ بحيث
    $b <_2 b_1$؛
  \item إذا كان $a_n <_1 a$، فليكن $b \in \Domain{M_2}$ بحيث
    $b_n <_2 b$؛
 \item إذا كان $a_i <_1 a <_1 a_{i+1}$ لبعض $i$، فليكن
   $b \in \Domain{M_2}$ بحيث $b_i <_2 b <_2 b_{i+1}$.
  \end{enumerate}
  ويمكن دائمًا إيجاد $b$ بالخاصية المطلوبة، لأن $\Struct{M_2}$ ترتيب
  خطي كثيف بلا طرفين. عرّف $q = p \cup \{ \langle a, b \rangle \}$،
  فتكون $q \in \PIso{I}$ التوسيع المطلوب لـ$p$. وهذا يثبت خاصية
  الذهاب. وخاصية الإياب مماثلة. ومن ثم
  $\Struct{M_1} \iso[p] \Struct{M_2}$؛ وبحسب
  \olref[pis]{thm:p-isom1}، $\Struct{M_1} \iso \Struct{M_2}$.
\end{proof}

\begin{prob}
  أتم برهان \olref[mod][bas][dlo]{thm:cantorQ} بالتحقق من أن
  $\PIso{I}$ تحقق خاصية الإياب.
\end{prob}

\begin{rem}
  لتكن $\Struct{S}$ أي ترتيب خطي كثيف !!{enumerable} بلا طرفين. عندئذ
  (بحسب \olref{thm:cantorQ}) يكون $\Struct{S} \iso \Struct{Q}$، حيث
  $\Struct{Q} = (\Rat, <)$ هو الترتيب الخطي الكثيف !!{enumerable} الذي
  مجاله مجموعة الأعداد النسبية~$\Rat$. تأمل الآن مرة أخرى
  الـ!!{structure}~$\Struct{R} = (\Real, <)$ من
  \olref[thm]{remark:R}. رأينا أن ثمة !!a{enumerable}
  !!{structure}~$\Struct{S}$ بحيث $\Struct{R} \elemequiv \Struct{S}$.
  لكن $\Struct{S}$ ترتيب خطي كثيف !!a{enumerable} بلا طرفين، ومن ثم فهو
  متشاكل (وبالتالي مكافئ أوليًا) مع الـ!!{structure}~$\Struct{Q}$.
  وبتعدي التكافؤ الأولي، $\Struct{R} \elemequiv \Struct{Q}$. (وكان
  بوسعنا إثبات ذلك مباشرة بإقامة $\Struct{R} \iso[p] \Struct{Q}$ بحجة
  الذهاب والإياب نفسها.)
\end{rem}
\end{document}
